\color{unidarkgreen}{\subsection{Автоматическое повторное включение}}
\color{black}

\subsubsection{Общие сведения}

\begin{enumerate}[label=\arabic{section}.\arabic{subsection}.\arabic{subsubsection}.\arabic*,  align=left, labelsep=4pt, leftmargin=0pt, itemindent=57pt]

\item
Автоматическое повторное включение (АПВ) предназначено для быстрого восстановления
первоначального состояния электрической сети после аварийного отключения путём автоматического
включения выключателя.

\item
Выбор действующих режимов АПВ для выключателя В1(ЛВ) ВН реализуется тремя входными
сигналами ФСУ:
\begin{itemize}
\item \textit{«АПВ В1(ЛВ) ВН ОНТ + ННШ»} (включение с контролем отсутствия напряжения на трансформаторе и наличия напряжения на шинах) – АПВ трансформатора;
\item \textit{«АПВ В1(ЛВ) ВН ОНШ» + ННТ»} (включение с контролем отсутствия напряжения на шинах и наличия напряжения на трансформаторе) – АПВ шин;
\item \textit{«АПВ В1(ЛВ) ВН без контролей»} (включение без контроля напряжений).
\end{itemize}

\item
Возможные режимы АПВ приведены в таблице \ref{desc:tbl1}.

\small
\renewcommand{\arraystretch}{1}
\setlength{\extrarowheight}{2pt}

\begin{longtable}{|>{\centering\arraybackslash}m{4.9cm}|>{\centering\arraybackslash}m{3.7cm}|>{\centering\arraybackslash}m{3.7cm}|>{\centering\arraybackslash}m{3.7cm}|}
\caption{Условия включения\hfill\vspace{-0.5\baselineskip}}\label{desc:tbl1}\\
\hline
\rowcolor{gray!30}
& \multicolumn{3}{c|}{Входные сигналы ФСУ} \\
\hhline{>{\arrayrulecolor{gray!30}}->{\arrayrulecolor{black}}|-|-|-|}
\rowcolor{gray!30}
\raisebox{3pt}{\makecell[c]{Режим АПВ}} & 
\makecell[c]{АПВ В1(ЛВ) ВН \\ ОНТ+ННШ} & 
\makecell[c]{АПВ В1(ЛВ) ВН \\ ОНШ+ННТ} & 
\makecell[c]{АПВ В1(ЛВ) ВН \\ без контролей} \\
\hline
\endfirsthead

\caption*{\hspace{3pt}\emph{Продолжение таблицы \ref{desc:tbl1}}\hfill\ }\vspace{-0.5\baselineskip} \\
\hline
\rowcolor{gray!30}
& \multicolumn{3}{c|}{Входные сигналы ФСУ} \\
\hhline{>{\arrayrulecolor{gray!30}}->{\arrayrulecolor{black}}|-|-|-|}
\rowcolor{gray!30}
\raisebox{3pt}{\makecell[c]{Режим АПВ}} & 
\makecell[c]{АПВ В1(ЛВ) ВН \\ ОНТ+ННШ} & 
\makecell[c]{АПВ В1(ЛВ) ВН \\ ОНШ+ННТ} & 
\makecell[c]{АПВ В1(ЛВ) ВН \\ без контролей} \\
\hline
\endhead

\hline
\makecell[c]{КС(УС)} & 0 & 0 & 0 \\
\hline
\makecell[c]{ОНТ+ННШ / КС(УС)} & 1 & 0 & 0 \\
\hline
\makecell[c]{ОНШ+ННТ / КС(УС)} & 0 & 1 & 0 \\
\hline
\makecell[c]{ОНТ+ННШ / ОНШ+ННТ /  КС(УС)} & 1 & 1 & 0 \\
\hline
\makecell[c]{Без контролей} & 0/1 & 0/1 & 1 \\
\hline
\end{longtable}
\normalsize

\item
Как видно из таблицы \ref{desc:tbl1} выбор режима АПВ, режим АПВ с КС(УС) введен всегда и выводится
посредством входного сигнала ФСУ \textit{«АПВ В1(ЛВ) ВН без контролей»}. А режимы ОНТ+ННШ и ОНШ+ННТ
могут быть введены одновременно. Режим без контролей применяется только для АПВ трансформатора.
Условия включения в различных режимах более подробно описаны в пункте {\color{red} 2.30} настоящего руководства
по эксплуатации.

\item
Пуск АПВ трансформатора осуществляется по цепи несоответствия при аварийном отключении
выключателя (по наличию входного сигнала ФСУ \textit{«Сигнал несоотв. В1(ЛВ) ВН»}). Срабатывание
происходит при набранной выдержке времени готовности АПВ с контролем заданного режима и
выполнения условий включения. Выдержка времени готовности АПВ после оперативного включения и
первого цикла регулируется уставкой \textbf{«Тгот АПВ В1(ЛВ) ВН»}, после второго цикла – \textbf{«Твосст гот.
АПВ В1(ЛВ) ВН»}.

\item
АПВ трансформатора может быть выполнено двукратным, ввод второго цикла осуществляется
программным ключом \textbf{«2-ой цикл АПВт В1(ЛВ) ВН»}. Первый цикл осуществляется при выполнении
условий включения с выдержкой времени \textbf{«Тср АПВт-1 В1(ЛВ) ВН»}. В случае сброса сигнала пуска АПВ
при отсчете времени \textbf{«Тср АПВт-1 В1(ЛВ) ВН»} 
цепь логики сбрасывается, и схема АПВ возвращается в исходное состояние (готовность АПВ с
набранной выдержкой готовности). Если первый цикл АПВ был успешным, то начинается набор
выдержки времени готовности к повторному действию.

\item
Если после первого цикла АПВ за время \textbf{«Тгот АПВ В1(ЛВ) ВН»} (выдержка времени к повторному
действию не успела набраться) повторно поступает сигнал пуска АПВ, то осуществляется 2-ой цикл с
выдержкой времени \textbf{«Тср АПВт-2 В1(ЛВ) ВН»}. В случае сброса сигнала пуска АПВ при отсчете времени
\textbf{«Тср АПВт-2 В1(ЛВ) ВН»} цепь логики сбрасывается, и схема возвращается в режим набора выдержки
времени готовности к повторному действию. Если второй цикл АПВ был успешным, то начинается набор
выдержки времени готовности к повторному действию.

\item
Второй цикл АПВ может быть выведен от внешнего сигнала \textit{«Вывод АПВт-2 В1(ЛВ) ВН»}. 

\item
Длительность импульса включения выключателя В1(ЛВ) ВН от АПВ трансформатор задается уставкой \textbf{«Тимп АПВт В1(ЛВ) ВН»}.

\begin{figure}[H]
\centering
\includegraphics[width=0.8\textwidth,height=0.8\textheight,keepaspectratio]{\fbpath/761.3011 АПВ/img/APV_T.pdf}
\captionof{figure}{Функциональная схема АПВ Т В1(ЛВ) ВН}\label{pic:APV_T}
\end{figure}

\item
Пуск АПВ шин осуществляется от внешнего сигнала о срабатывании ДЗШ (входной сигнал ФСУ
\textit{«Пуск АПВ В1(ЛВ) ВН от ДЗШ»}). При этом АПВ трансформатор запрещается. АПВ шин выполнено
однократным и осуществляется при выполнении условий включения с выдержкой времени \textbf{«Тср АПВш
В1(ЛВ) ВН»}. Длительность импульса включения выключателя В1(ЛВ) ВН от АПВ шин задается уставкой
\textbf{«Тимп АПВш В1(ЛВ) ВН»}.

\begin{figure}[H]
\centering
\includegraphics[width=0.8\textwidth,height=0.8\textheight,keepaspectratio]{\fbpath/761.3011 АПВ/img/APV_SSh.pdf}
\captionof{figure}{Функциональная схема АПВ СШ В1(ЛВ) ВН}\label{pic:APV_SSh}
\end{figure}

\item
В устройстве предусмотрено ограничение ожидания условий включения, которое вводится
программным ключом \textbf{«Огр. ожид. усл. вкл. В1(ЛВ) ВН»}. Если в течение выдержки времени \textbf{«Тожид
усл. вкл. В1(ЛВ) ВН»} после аварийного отключения не возникают условия включения, то действие АПВ
запрещается.

\item
Предусмотрена возможность запрета АПВ при длительно отключенном выключателе, который
вводится программным ключом \textbf{«Огр. ожид. пуска АПВ В1(ЛВ) ВН»}. Действие на запрет АПВ
осуществляется с выдержкой времени \textbf{«Тожид усл. пуска АПВ В1(ЛВ) ВН»}.


\small
\begin{longtable}{|>{\centering\arraybackslash}m{5.3cm}|>{\centering\arraybackslash}m{4.2cm}|>{\centering\arraybackslash}m{1cm}|>{\centering\arraybackslash}m{3.3cm}|>{\centering\arraybackslash}m{1.8cm}|}
% Первая страница: полный заголовок
\caption{Перечень уставочных параметров АПВ\hfill\vspace{-0.5\baselineskip}}\label{desc:tbl1}\\
\hline
\rowcolor{gray!30}
Наименование уставки & Диапазон & Шаг & Значение по умолчанию & Единица измерения \\
\hline
\endfirsthead

% Промежуточные страницы: «Продолжение таблицы»
\caption*{\hspace{3pt}\emph{Продолжение таблицы \ref{desc:tbl1}}\hfill\ }\vspace{-0.5\baselineskip} \\
\hline
\rowcolor{gray!30}
Наименование уставки & Диапазон & Шаг & Значение по умолчанию & Единица измерения \\
\hline
\endhead

%===m>ДТЗ|-
%===m
\end{longtable}
\normalsize	

\end{enumerate}